\documentclass[UTF8,11pt]{ctexart}
\usepackage{amsmath,amssymb,mathtools,geometry}
\geometry{a4paper,margin=20mm}\allowdisplaybreaks
\begin{document}
\section*{中继充电站无人机网络设计完整线性弧模型}

\paragraph{来源} PDF第6--9页，§4，式(8)--(20)。目标采用(18)--(20)的乘积线性化。
\subsection*{集合、参数与变量}
$S$ 为既有仓/枢纽，$I$ 为候选充电站，$N$ 为需求地点，$S\cap I=\varnothing$。
$\Gamma=S\cup I$。$E$ 为可行有向飞行弧：枢纽至候选站、候选站之间，
每段长度 $l_e\le2R$，并已删除障碍物不允许的弧；不含进入枢纽或自环的弧。
$c_{in}=1$ 当候选站 $i$ 至需求点 $n$ 的距离不超过 $R$，否则为零。
$\theta\in[0,1]$ 为权重，$\beta_1,\beta_2>0$ 为给定的规范化常数。
$\delta^+(k),\delta^-(k)$ 分别表示以 $k$ 为起/终点的弧。
$x_k$ 表示候选站启用；$u_{si}$ 表示终端站 $i$ 由枢纽 $s$ 连接；
$y_e^{si}$ 表示该连接路径使用弧 $e$，$z_e^{si}=y_e^{si}u_{si}$。
这些变量均为二元。
\subsection*{完整模型}
\begin{align}
\min\quad&
\frac{\theta}{\beta_1}\sum_{s\in S,i\in I,e\in E}l_e z_e^{si}
+\frac{1-\theta}{\beta_2}\sum_{k\in I}x_k,\\
\sum_{e\in\delta^+(k)}y_e^{si}-\sum_{e\in\delta^-(k)}y_e^{si}
&=\begin{cases}u_{si},&k=s,\\-u_{si},&k=i,\\0,&k\in\Gamma\setminus\{s,i\},\end{cases}
&&s\in S,\ i\in I,\ k\in\Gamma,\\
\sum_{e\in\delta^+(s)}y_e^{si}&\le1 &&s\in S,\ i\in I,\\
y_e^{si}&\le u_{si} &&s\in S,\ i\in I,\ e\in E,\\
\sum_{s\in S}u_{si}&\le1 &&i\in I,\\
\sum_{s\in S,i\in I}c_{in}u_{si}&\ge1 &&n\in N,\\
y_e^{si}&\le x_k &&s\in S,\ i\in I,\ k\in I,\ e\in\delta^-(k),\\
z_e^{si}&\le y_e^{si},\qquad z_e^{si}\le u_{si}
&&s\in S,\ i\in I,\ e\in E,\\
z_e^{si}&\ge y_e^{si}+u_{si}-1 &&s\in S,\ i\in I,\ e\in E,\\
x_k,u_{si},y_e^{si},z_e^{si}&\in\{0,1\}. &&
\end{align}
可行域由上述完整弧集合和约束定义。未访问候选站可关闭；所有实际中继或终端节点必须启用。
\paragraph{原文索引澄清}
原式(14)印作 $\sum_{s,e}y_e^{si}\le Mx_i$，只用终端索引 $i$，但第9页紧接的说明明确要求
路径进入的每一个候选节点启用。上式将这个逐节点关系展开为 $y_e^{si}\le x_k$，避免遗漏中继站。
流守恒在所有枢纽/候选节点上显式展开，禁止一个commodity从其他枢纽偷取流。
索引按上述网络节点和商品集合展开。

\end{document}
